% ============================================================================
% Fichier  : partie4/P04_C11_droits_personnes.tex
% Titre    : Auditer les Droits des Personnes
% Auteur   : MINKA MI NGUIDJOÏ Thierry Emmanuel
% Version  : 1.0 -- novembre 2026
% Corrigés : corriges/Corr_C11.tex
% ============================================================================
% Sources : loi n° 2024/017, art. 21, 37 à 47, 56, 62, 64 (texte officiel) ;
% ISO/IEC 27701 (traitement des demandes, paraphrasé).
% Faits nouveaux (bible) : 31 demandes du 23/06 au 31/08/2026, 12 traitées,
% délai médian 26 jours ; demande test de KORO AUDIT, 27 juillet 2026.
% ============================================================================

\chapter{Auditer les Droits des Personnes}
\label{chap:droits}
\chaptermark{Les droits des personnes}

\epigraph{%
    « Un droit qu'on ne peut pas exercer\\
    n'est qu'une phrase dans une loi. »%
}{Principe du cours}

\niveauChapitre{Intermédiaire}{%
    Chapitres~\ref{chap:principes} et~\ref{chap:audit}
}

\begin{boxobjectifs}
\begin{itemize}
    \item Énoncer les droits des articles~37 à~46 et leurs limites.
    \item Vérifier qu'un canal de demande existe, qu'il est connu et que
          l'identité du demandeur est contrôlée sans excès.
    \item Mesurer les délais de réponse en l'absence de délai réglementaire,
          avec un critère de substitution déclaré.
    \item Tester le traitement d'une demande de bout en bout.
    \item Relier les défaillances aux sanctions administratives et pénales
          encourues.
\end{itemize}
\end{boxobjectifs}

\bigskip

% ============================================================================
\section*{Avant de commencer : trente et une lettres}
% ============================================================================

Bafoussam, lundi 27~juillet~2026. La lettre de M.~TAGNE n'est pas restée
seule. L'audience de février a été relayée par la radio locale, et d'autres
clients ont écrit. Me~ETOUNDI a ouvert un classeur~: entre le 23~juin et le
27~juillet, vingt-deux lettres, courriels et messages vocaux laissés au
standard. Certaines demandent l'accès, d'autres l'effacement, deux
s'opposent aux appels de prospection, une demande la rectification d'un
nom mal orthographié qui bloque un virement.

Il n'existe ni procédure, ni délai, ni modèle de réponse. Les demandes
arrivent par six canaux différents~; trois ont été perdues entre l'agence
et le siège. La seule qui ait reçu une réponse complète est celle d'un
client qui connaissait personnellement le directeur général.

\begin{boxalert}
\textbf{Violation typique.} Considérer que l'organisation respecte les droits
des personnes parce qu'elle n'a « jamais refusé » une demande.

\smallskip
\textbf{Règle~:} le respect des droits se prouve par une procédure écrite,
un registre des demandes, des réponses datées et un délai mesuré. Une
demande perdue est un refus de fait~; une demande traitée par faveur
n'est pas un droit exercé.
\end{boxalert}

% ============================================================================
\section{Les Droits de la Personne Concernée}
\label{sec:c06-droits}
% ============================================================================

Le titre~III de la loi (articles~37 à~47) organise les droits. L'article~21
impose d'en informer la personne dès la collecte~; l'auditeur vérifie donc
à la fois l'information et l'exercice.

\begin{tableaucours}{L{3.0cm}cL{8.0cm}}{%
    \tableauHeader{Droit} &
    \tableauHeader{Article} &
    \tableauHeader{Contenu et limites}
}
Accès & 39 & Confirmation du traitement, communication des données et de
leur origine, finalités, destinataires~; copie contre paiement limité au
coût de reproduction \\
Rectification & 42 & Rectifier, compléter, mettre à jour, verrouiller ou
supprimer les données inexactes ou incomplètes, sur justification de
l'identité \\
Effacement & 37, 38 & Cessation de la diffusion et effacement si les
données ne sont plus nécessaires, si le consentement est vicié, retiré ou
expiré, ou sans base légale~; sous réserve des obligations de conservation \\
Opposition & 40 & À tout moment, pour les catégories particulières de
données, pour la prospection, et pour tout motif prévu par la loi~; droit
propre des parents d'un mineur \\
Prospection & 41 & Communication à des tiers pour la prospection soumise à
consentement préalable \\
Portabilité & 43 & Recevoir gratuitement ses données dans un format
structuré, et les faire transmettre à un autre responsable~; exclue pour
les missions d'intérêt public \\
Décision automatisée & 44 & S'opposer à une décision fondée exclusivement
sur un traitement automatisé (Chapitre~\ref{chap:profilage}) \\
Décès & 45 & Fin du traitement au décès, sauf exceptions~; directives
post-mortem~; mise à jour par les ayants droit \\
Limitation & 46 & Geler le traitement quand l'exactitude ou la finalité est
contestée \\
\end{tableaucours}
\vspace{-0.8em}
\begin{center}{\small\itshape Tableau~11.1. Les droits de la personne concernée}\end{center}

\begin{boxjuris}
\textbf{\loicam{n°~2024/017}, article~43}~\cite{loi2024017}~:
\textit{« (1) La personne concernée a le droit de recevoir gratuitement les
données à caractère personnel la concernant qu'elle a fournies à un
responsable du traitement, dans un format structuré, couramment utilisé et
lisible par machine. (2) La personne concernée a le droit à la portabilité
de ses données. (3) Lorsque la personne concernée exerce son droit à la
portabilité des données en application de l'alinéa 1 ci-dessus, elle a le
droit d'obtenir que les données à caractère personnel soient transmises
directement d'un responsable du traitement à un autre, si cela est
techniquement possible. [\ldots] »}

\smallskip
\textbf{Portée.} L'accès peut être facturé au coût de reproduction
(article~39)~; la portabilité est gratuite. Pour une microfinance, la
portabilité permet à un client de faire transmettre son historique à un
autre établissement~: l'organisation doit être capable de l'exporter.
\end{boxjuris}

L'article~45 contient une disposition originale~: le traitement des données
d'une personne « cesse dès constatation de son décès », sauf obligation
légale, défense en justice ou directives que la personne a laissées, et les
ayants droit peuvent demander la mise à jour des informations du défunt.
Pour un établissement financier, cela suppose une procédure de succession
qui distingue ce que la loi impose de garder et ce qui doit cesser.

% ============================================================================
\section{Le Délai, un Critère à Déclarer}
\label{sec:c06-delai}
% ============================================================================

L'article~47 renvoie à la voie réglementaire « les délais et les modalités
d'examen des demandes ». Tant que ce texte n'est pas publié, la loi ne fixe
aucun délai. L'auditeur a deux possibilités, à choisir dans le plan
d'audit~:

\begin{itemize}
    \item coter \NEK{} de critère toute exigence de délai réglementaire, ce
          qui est exact mais ne dit rien de la qualité du service~;
    \item déclarer un \textbf{critère de substitution}, par exemple
          « réponse complète dans les 30 jours calendaires », en précisant
          qu'il s'agit d'une référence de bonne pratique et non d'une
          obligation légale, et mesurer les délais contre ce critère.
\end{itemize}

KORO AUDIT a retenu la seconde, en la déclarant dans le plan d'audit. Le
rapport dira donc « 30~jours (critère de l'audit) », jamais « 30~jours
(délai légal) ».

% ============================================================================
\section{Auditer le Traitement des Demandes}
\label{sec:c06-auditer}
% ============================================================================

\begin{boxprocedure}
\textbf{Auditer l'exercice des droits.}
\begin{enumerate}
    \item Vérifier l'information~: les droits et la manière de les exercer
          figurent-ils sur les formulaires et la notice (article~21)~?
    \item Recenser les canaux réels~: courrier, courriel, agence, téléphone,
          application, réseaux sociaux.
    \item Obtenir le registre des demandes, ou le reconstituer à partir des
          canaux.
    \item Pour chaque demande, relever la date de réception, le type, la
          vérification d'identité, la date et le contenu de la réponse.
    \item Calculer les délais, la part des demandes traitées et la part
          traitées dans le délai retenu.
    \item Réaliser une demande test, prévue au mandat, et la suivre jusqu'à
          la réponse.
    \item Vérifier que les effacements et les oppositions ont été exécutés
          dans les systèmes, et pas seulement acceptés par lettre.
\end{enumerate}
\end{boxprocedure}

La vérification d'identité mérite attention. L'article~42 vise la personne
« justifiant de son identité »~: l'organisation doit s'assurer qu'elle ne
communique pas les données d'un client à un imposteur, sans exiger plus que
nécessaire. Demander une copie de la carte nationale d'identité pour une
simple opposition aux SMS est excessif~; la demander pour communiquer un
relevé de compte est justifié.

Le dernier point de la procédure est le plus souvent oublié. Une lettre
« nous avons bien pris en compte votre demande d'effacement » ne prouve pas
l'effacement~: l'auditeur cherche le client dans la base, dans les
sauvegardes de moins d'un mois, dans le fichier du prestataire.

% ============================================================================
\section{Les Demandes de NYANGA}
\label{sec:c06-nyanga}
% ============================================================================

\begin{boxscenario}{Mission NYANGA. Le registre reconstitué, 31 août 2026}
\textbf{Constatations.}
\begin{itemize}
    \item 31~demandes reçues entre le 23~juin et le 31~août~: 14 d'accès,
          8 d'effacement, 5 d'opposition à la prospection, 3 de
          rectification, 1 de portabilité.
    \item 12 ont reçu une réponse complète, 7 une réponse d'attente, 12
          aucune réponse. Délai médian des réponses complètes~: 26~jours~;
          4 au-delà de 30~jours.
    \item Les 5 oppositions ont été acceptées par lettre~; 3 des 5 numéros
          figurent encore dans la liste d'appels de prospection du 28~août.
    \item La demande de portabilité a été refusée au motif que « le
          logiciel ne permet pas l'export ».
    \item La demande test de KORO AUDIT, envoyée le 27~juillet par courriel à
          l'adresse du service client, a reçu un accusé de réception
          automatique et aucune réponse au 31~août.
\end{itemize}

\textbf{Analyse} (confiance élevée). L'organisation n'a pas de procédure~;
les réponses dépendent de la personne qui reçoit la demande. Les
oppositions non exécutées exposent l'établissement à l'article~64, qui
punit le traitement maintenu malgré l'opposition à la prospection.
\end{boxscenario}

\begin{boxlecture}{Le domaine D2, droits des personnes}
\begin{tableaucours}{L{7.8cm}cL{3.2cm}}{%
    \tableauHeader{Exigence} &
    \tableauHeader{Cote} &
    \tableauHeader{Preuve}
}
Les personnes sont informées de leurs droits & \NON{} & Formulaire 2016 \\
Une procédure de traitement des demandes existe & \NON{} & Aucune \\
Un registre des demandes est tenu & \NON{} & Reconstitué par l'audit \\
Les demandes d'accès reçoivent une réponse & \NON{} & 12 sur 31 \\
Les réponses respectent le critère de 30 jours & \NON{} & 8 sur 12 \\
Les oppositions sont exécutées dans les systèmes & \NON{} & 3 numéros
encore appelés \\
La portabilité est possible & \NON{} & Refus du 14~août \\
L'identité est vérifiée de façon proportionnée & \NEK{} & Pratiques
variables, non documentées \\
Les effacements acceptés sont exécutés & \NEK{} & Sauvegardes non
examinées \\
La procédure de succession respecte l'article~45 & \OUI{} & Note de 2025,
appliquée sur 3 dossiers \\
\end{tableaucours}

\medskip
\textbf{Décompte.} Un \OUI{}, sept \NON{}, deux \NEK{} sur dix.
\hfill\lecture{0{,}10}{0{,}20}{1{,}16}

\medskip
\textbf{Lecture.} $U = 0{,}20$ atteint juste le seuil~: le domaine est
opaque selon la règle, mais de peu. La non-conformité est établie sur sept
exigences, dont une, l'opposition non exécutée, relève d'une sanction
pénale. La recommandation prioritaire ne dépend pas des deux \NEK{}.
\end{boxlecture}

\begin{boxjuris}
\textbf{\loicam{n°~2024/017}, article~56}~\cite{loi2024017}~:
\textit{« Est passible d'une amende d'un million (1~000~000) à dix millions
(10~000~000) de francs CFA, tout responsable de traitement ou sous-traitant
qui refuse de mettre à la disposition de la personne concernée, les
informations sollicitées la concernant. »}

\smallskip
\textbf{Portée.} Une demande d'accès laissée sans réponse peut être
qualifiée de refus. L'article~62 ouvre en outre à la personne la voie du
référé « en cas d'atteinte grave aux droits », et la réparation devant la
juridiction compétente~: c'est la voie qu'annonçait M.~TAGNE.
\end{boxjuris}

% ============================================================================
\section{Ce que les Droits demandent à chaque Métier}
\label{sec:c06-metiers}
% ============================================================================

\begin{boxmetier}{Délégué à la protection des données}
Centraliser les demandes sur un canal connu, tenir le registre, fixer un
délai cible, vérifier l'exécution dans les systèmes, et rendre compte
chaque mois du nombre de demandes et des délais.
\end{boxmetier}

\begin{boxmetier}{Juriste conformité}
Préparer des modèles de réponse par type de droit, et une grille des
données à conserver malgré une demande d'effacement (comptabilité,
procédure judiciaire, lutte contre le blanchiment).
\end{boxmetier}

\begin{boxmetier}{Auditeur de la protection des données}
Toujours vérifier l'exécution, pas seulement l'acceptation~; et déclarer
dans le plan d'audit tout critère de délai qui n'est pas légal.
\end{boxmetier}

\begin{boxresume}
\begin{itemize}
    \item La loi reconnaît l'accès, la rectification, l'effacement,
          l'opposition, la portabilité gratuite, la limitation, des droits
          face aux décisions automatisées et des règles en cas de décès.
    \item Les délais de réponse relèvent d'un texte d'application~; en son
          absence, l'auditeur déclare un critère de substitution.
    \item Un droit exercé se prouve par un registre, une réponse datée et
          l'exécution dans les systèmes.
    \item Le refus de communiquer les informations (article~56) et le
          traitement maintenu malgré l'opposition (article~64) sont
          sanctionnés.
    \item À NYANGA, 12 demandes sur 31 ont reçu une réponse complète, et des
          oppositions acceptées n'ont pas été exécutées.
\end{itemize}
\end{boxresume}

\begin{boxinfo}
\textbf{Points de contrôle.} Vous avez compris ce chapitre si vous savez~:
\begin{itemize}
    \item citer les limites du droit à l'effacement~;
    \item expliquer pourquoi un délai de 30~jours ne peut pas être présenté
          comme légal~;
    \item vérifier l'exécution d'une opposition dans les systèmes.
\end{itemize}
\end{boxinfo}

\begin{boxexo}[title={\ding{45}\quad Exercice~11.1. Quel droit \niveaubase}]
Quel droit chaque demande exerce-t-elle~? (a) « Donnez-moi la liste des
personnes à qui vous avez transmis mon numéro. »~; (b) « Je ne veux plus de
vos SMS. »~; (c) « Transmettez mon historique de crédit à ma nouvelle
banque. »~; (d) « Mon père est décédé~; supprimez son compte de vos
campagnes. »~; (e) « Mon salaire est faux dans votre dossier. »
\end{boxexo}

\begin{boxexo}[title={\ding{45}\quad Exercice~11.2. Mesurer les délais \niveaumoyen}]
Voici les délais de réponse complète, en jours, de dix demandes~: 5, 12,
18, 21, 26, 26, 31, 38, 44, 60. Calculez le délai médian et la part des
réponses dans le critère de 30~jours. Rédigez la phrase du rapport.
\end{boxexo}

\begin{boxexo}[title={\ding{45}\quad Exercice~11.3. L'effacement de M. TAGNE \niveaumoyen}]
Dressez la liste des données de M.~TAGNE que NYANGA peut effacer et de
celles qu'elle doit conserver, avec pour chacune le texte qui le justifie
ou la mention \NEK{} si vous ne connaissez pas le texte.
\end{boxexo}

\begin{boxexo}[title={\ding{45}\quad Exercice~11.4. La demande test \niveauexpert}]
Rédigez la clause du mandat qui autorise la demande test, puis le papier de
travail qui en rend compte. Quelles précautions prenez-vous pour que le
test ne soit pas lui-même un traitement déloyal~?
\end{boxexo}

\bigskip

\begin{boxinfo}
\textbf{Pour aller plus loin.}
\begin{itemize}
    \item L'ISO/IEC~27701~\cite{iso27701}, pour l'organisation du traitement
          des demandes des personnes.
\end{itemize}
\end{boxinfo}

\begin{boxinfo}
\textbf{Transition.} Les droits protègent la personne contre l'organisation.
La sécurité la protège contre tous les autres~: intrus, agents curieux,
prestataires négligents. C'est l'objet du Chapitre~\ref{chap:securite}.
\end{boxinfo}
